\scanpage{43}
\mspage{9}
élémentaire de $S^2$ \inserted{de (23)}, on trouve donc
que la donnée d'une gerbe sur $S^2$ équivaut
à celle d'un \uline{automorphisme} de
l'extension triviale \deleted{\unsure{sur}} de $\mathbb Z$ par $A$,
donc à un élément de $\operatorname{Hom}(\mathbb Z,A)\simeq A$,
d'où la valeur correcte
\[
(25)\qquad H^2(S^2,\deleted{\mathbb Z}\inserted{A})\simeq A,
\]
magnifique succès ! Sans aller plus
loin dans des vérifications, ceci me
convainc que l'intuition qu'il y
a une théorie des \deleted{limites inductives}\inserted{$\underset{\longrightarrow}{\mathrm{Lim}}$}
des types d'homotopie (avec un
théorème à la \uline{van Kampen}, disant
que le type d'homotopie d'une
$\unsure{\underline{\mathrm{Lim}}|\underset{\longrightarrow}{\mathrm{Lim}}}$ de Topos (Topos ``recollé'') est
la $\unsure{\underline{\mathrm{Lim}}|\underset{\longrightarrow}{\mathrm{Lim}}}$ des types d'homotopie),
est bien correcte. En principe, pour
\mspage{42}
